\begin{table}[!h]
\centering\centering\centering
\caption{Effect on Test Scores}
\centering
\begin{threeparttable}
\fontsize{10}{12}\selectfont
\begin{tabular}[t]{lcccccc}
\toprule
\multicolumn{1}{c}{ } & \multicolumn{6}{c}{Experiment (Liberia)} \\
\cmidrule(l{3pt}r{3pt}){2-7}
  & (1) & (2) & (3) & (4) & (5) & (6)\\
\midrule
Peer Baseline Score & -0.052 &  &  &  & -0.007 & -0.025\\
 & (0.084) &  &  &  & (0.196) & (0.201)\\
Level &  & -0.042 & -0.044 &  & -0.048 & -0.045\\
 &  & (0.105) & (0.104) &  & (0.144) & (0.144)\\
Level $\times$ Std. Baseline Score &  &  & -0.030 &  &  & -0.040\\
 &  &  & (0.048) &  &  & (0.041)\\
Class Size &  &  &  & -0.003 & -0.001 & -0.001\\
 &  &  &  & (0.005) & (0.005) & (0.005)\\
\hline
\midrule
Number of Students & 2785 & 2785 & 2785 & 2785 & 2785 & 2785\\
Number of Schools & 54 & 54 & 54 & 54 & 54 & 54\\
Number of Test Scores & 1 & 1 & 1 & 1 & 1 & 1\\
\bottomrule
\end{tabular}
\begin{tablenotes}[para]
\item Notes: All specifications control for individual predicted endline test score, 
           grade dummies, and randomization strata fixed effects. Each specification also controls for the    
           expected value of the respective independent variable used in the specification viz. peer predicted              score, level and class size. Standard errors are clustered at the academy-grade group level. ***, **,            and * indicate significance at 1\%, 5\%, and 10\%.
\end{tablenotes}
\end{threeparttable}
\end{table}
